\begin{textsample}{NEG Dim 1 – 2012\_c6 – Score -3.00 – 2012\_c6/enade\_2012\_admin\_bb7f435184.txt}  \label{ex:f1_neg_254}
CAPÍTULO 3 ANÁLISE TÉCNICA DA PROVA Este capítulo tem por objetivo apresentar o desempenho dos estudantes concluintes de Administração no ENADE/2012.
Para isso, foram calculadas as estatísticas básicas da prova em seu todo, bem como as estatísticas dos componentes relacionadas à Formação Geral, ao de Conhecimento Específico da Área e das questões discursivas isoladamente.
Nas tabelas, são apresentados o tamanho da população inscrita e de presentes, e as seguintes estatísticas das notas9 : média do desempenho na prova, erro padrão da média, desvio padrão, nota mínima, mediana e nota máxima.
Tais estatísticas contemplam o total de estudantes concluintes da Área de Administração inscritos e presentes à prova do ENADE 2012, tendo -se em vista as seguintes agregações : ( a ) as Grandes Regiões e o país como um todo ; ( b ) a Categoria Administrativa ; e ( c ) a Organização Acadêmica.
Em relação aos gráficos de distribuição de notas, o intervalo considerado foi de 10 unidades, aberto à esquerda e fechado à direita, com exceção do primeiro intervalo, [ 0 ; 10 ], fechado em ambos os extremos.
Para os gráficos de distribuição das notas das questões discursivas, foram consideradas mais duas categorias : questão em branco e nota zero. 3.1 ESTATÍSTICAS BÁSICAS DA PROVA 3.1.1 Estatísticas Básicas Gerais A Tabela 3.1 apresenta as estatísticas básicas da prova por grande Região.
A população total de inscritos foi de 176.254.
Destes, 140.180 \textbf{estiveram} presentes, sendo 20,5 \% o índice de não comparecimento.
A Região de maior abstenção foi a Norte ( 41,3 \% ) e a de menor abstenção foi a Sul ( 16,9 \% ).
A média das notas da prova como um todo ( nas seções seguintes serão analisados os componentes de Formação Geral e de Conhecimento Específico ) foi 34,6, sendo que os alunos da região Centro-Oeste obtiveram a média mais baixa ( 32,4 ) e os da região Nordeste obtiveram a média mais alta ( 35,4 ).
As demais médias foram : 35,2 na região Sul, 34,8 na região Sudeste e 32,7 na região Norte.
O desvio padrão para o Brasil como um todo foi 13,6, sendo o maior desvio padrão encontrado nas regiões Nordeste e Sudeste ( 13,8 ) e o menor na região Norte ( 11,8 ), indicando uma menor dispersão das notas desta última região. 9 Essas estatísticas e outras \textbf{estão} definidas no Capítulo 1.
A região que obteve a maior nota máxima foi a Nordeste ( 89,3 ), ao passo que a região que atingiu a menor nota máxima foi a Norte ( 82,6 ).
A mediana do Brasil como um todo foi 33,0, sendo a maior mediana obtida na região Nordeste ( 34,1 ) e a menor obtida na Centro-Oeste ( 30,7 ).
A nota mínima foi zero em todas as regiões, sem exceção.
O comportamento das notas dos estudantes de todo o Brasil pode ser observado no Gráfico 3.1 que apresenta um histograma com a distribuição das mesmas.
Essa é uma distribuição unimodal com moda no intervalo ( 20 ; 30 ].
Destaca -se, também, o intervalo ( 30 ; 40 ] com frequência relativa muito próxima à modal.
O coeficiente de assimetria da distribuição das notas é positivo ( 0,49 ) e não podemos considerar que a distribuição é aproximadamente simétrica.
As distribuições por Grande Região também apresentam assimetria positiva, concentração pouco maior do lado esquerdo do histograma e mais espalhada do lado direito.
Os Gráficos 3.2, 3.3 e 3.4 apresentam informações referentes à média da nota final dos participantes, desagregando os resultados de acordo com, respectivamente, as Grandes Regiões do país, a Categoria Administrativa e a Organização Acadêmica.
Os gráficos apresentam o valor da média das notas como uma barra e os extremos do intervalo de confiança de 95 \% como linhas verticais unidas por uma linha horizontal na forma da letra H maiúscula.
Considerando -se o gráfico de notas segundo Grande Região ( Gráfico 3.2 ), observa -se que existe diferença estatisticamente significativa ao nível de 95 \% entre a maior média, obtida na região Nordeste ( 35,4 ) e a menor, obtida na região Centro-Oeste ( 32,4 ).
Levando -se em conta os estudantes por Categorias Administrativas ( Gráfico 3.3 ), observa -se que existe diferença estatisticamente significante entre as médias das notas das IES Públicas e Privadas.
Pode ser observado ainda que a média dos alunos de IES Públicas ( 40,8 ) é maior que a dos alunos de IES Privadas ( 33,9 ).
A diferença entre as médias das regiões Nordeste e Centro-Oeste ( 3,0 ), a maior e a menor média, é inferior a diferença entre IES Públicas e Privadas ( 6,9 ).
Tendo como foco o Gráfico 3.4, que apresenta as notas médias das provas segundo Organização Acadêmica, constata -se que a diferença da média das notas dos estudantes provenientes de Universidades é estatisticamente significativa ao nível de 95 \% das médias das notas dos estudantes provenientes de Centros Universitários e Faculdades.
Já a diferença entre a média dos estudantes de Centros Universitários e Faculdades não é significativa estatisticamente.
A maior média foi obtida pelos estudantes de Universidades ( 35,5 ), e a menor, pelos de Faculdades ( 33,9 ). 3.1.2 Estatísticas Básicas no Componente de Formação Geral A Tabela 3.2 apresenta as estatísticas básicas em relação ao componente da prova que avalia a Formação Geral dos estudantes concluintes.
Os alunos de todo Brasil obtiveram desempenho médio de 42,9.
Quanto à variabilidade, o desvio padrão das notas dos estudantes do Brasil como um todo foi 16,2.
A maior média foi obtida na região Nordeste ( 43,5 ), e a menor, na região Centro-Oeste ( 41,5 ).
As demais médias foram : 43,1 na região Sul, 42,9 na região Sudeste e 42,2 na região Norte.
Já o maior desvio padrão foi obtido na região Nordeste ( 16,5 ) e o menor na região Norte ( 15,3 ).
Os demais desvios padrões foram : 16,3 nas regiões Centro-Oeste e Sudeste, e 15,9 na região Sul.
A maior nota no componente de Formação Geral da prova do ENADE foi obtida por pelo menos um aluno da região Sudeste e um aluno da região Sul ( 95,0 ) enquanto que a menor nota máxima foi obtida na região Norte ( 92,0 ).
Nas outras regiões as notas máximas foram : 93,0 tanto para a região Nordeste quanto para a região Centro-Oeste.
A mediana do Brasil como um todo foi 43,5, sendo a menor mediana encontrada na região Centro-Oeste ( 42,0 ) e a maior encontrada na região Nordeste ( 44,5 ).
A nota mínima nesta parte foi zero em todas as regiões, sem exceção.
O Gráfico 3.5 propicia a avaliação do desempenho dos estudantes no componente de Formação Geral a partir do histograma da distribuição das notas correspondentes.
A distribuição é unimodal, com moda em ( 40 ; 50 ], enquanto na prova como um todo a moda foi alcançada no intervalo ( 20 ; 30 ].
Nota -se, ainda, que no gráfico 3.5 as notas apresentam uma maior dispersão do que no Gráfico 3.1 ( distribuição das notas da prova ), confirmado pela comparação dos desvios padrões : 13,6 para a nota da prova como um todo e 16,2 para o componente de Formação Geral.
Para o componente de Formação Geral, o coeficiente de assimetria da distribuição das notas dos estudantes é negativo ( – 0,14 ).
Por ser um coeficiente pequeno, a distribuição apresenta uma leve concentração à direita e cauda maior à esquerda, porém podemos considerar que a distribuição é aproximadamente simétrica.
Em todas as Grandes Regiões os histogramas também possuem assimetria levemente negativa ( entre -0,18 para o Nordeste e -0,09 para o Norte ).
Nos Gráficos 3.6, 3.7 e 3.8 são apresentadas as informações referentes ao desempenho dos Concluintes no componente de Formação Geral, em diferentes agregações : Grande Região do país, Categoria Administrativa e Organização Acadêmica.
Observa -se pelo Gráfico 3.6 que existe diferença estatisticamente significativa ao nível de 95 \% entre a maior média das notas na componente de Formação Geral, obtida na região Nordeste ( 43,5 ) e a menor, obtida na região Centro-Oeste ( 41,5 ).
Vemos que o intervalo de confiança mais largo é o da região Norte ; já o intervalo mais estreito é observado na região Sudeste.
Este fato \textbf{está} relacionado, também, com o tamanho da população envolvida, menor na região Norte do que na Sudeste.
No Gráfico 3.7, que representa as notas médias no Componente de Formação Geral segundo Categoria Administrativa do país, observa -se que existe diferença estatisticamente significativa entre as médias.
Assim como ocorreu para as notas da prova como um todo, para o componente de Formação Geral, os concluintes de Administração, das IES Públicas ( 47,0 ) obtiveram uma média maior do que os das IES Privadas ( 42,4 ).
Considerando -se o tipo de Organização Acadêmica, nota -se, no Gráfico 3.8, uma diferença estatisticamente significativa entre a maior e a menor média.
Nas Universidades, a média ( 43,5 ) é mais elevada do que em Centros Universitários ( 42,2 ) e Faculdades ( 42,5 ). 3.1.3 Estatísticas Básicas do Componente de Conhecimento Específico A Tabela 3.3 apresenta as estatísticas básicas referentes ao componente de Conhecimento Específico da área de Administração.
A média do desempenho dos alunos do Brasil como um todo foi 31,9.
A maior média foi obtida na região Nordeste ( 32,6 ), e a menor, na região Centro-Oeste ( 29,3 ).
As demais médias foram : 32,5 na região Sul, 32,0 na região Sudeste e 29,5 na região Norte.
Quanto à variabilidade das notas, o desvio padrão do Brasil como um todo foi 15,1, sendo o maior desvio padrão observado nas regiões Nordeste, Sudeste e Sul ( 15,3 ) e o menor na região Norte ( 13,1 ).
O desvio padrão da região Centro-Oeste foi ( 14,2 ).
A mediana das notas dos estudantes de todo o Brasil foi 29,3.
A maior mediana ocorreu na região Nordeste ( 30,8 ) e a menor na região Centro-Oeste ( 28,1 ).
As demais medianas foram 30,3 na região Sul, 29,3 na região Sudeste e 28,3 na região Norte.
A nota máxima do Brasil como um todo foi 95,0, sendo obtida por pelo menos um aluno da região Sul.
As demais notas máximas foram : 93,7 na região Sudeste, 92,2 na região Nordeste, 91,9 na região Centro-Oeste e 91,9 na região Norte.
A nota mínima foi zero em todas as regiões, sem exceção.
Assim como os Gráficos 3.1 e 3.5, o Gráfico 3.9, apresentado a seguir, proporciona uma avaliação do desempenho de concluintes em relação ao componente de Conhecimento Específico com um histograma da distribuição das notas correspondentes.
Esta também é uma distribuição unimodal, e o grupo modal é o ( 20 ; 30 ], abaixo do grupo modal para a Formação Geral.
O coeficiente de assimetria da distribuição das notas do componente de Conhecimento Específico é positivo ( 0,65 ).
Nota -se pelo histograma ( Gráfico 3.9 ) que esta distribuição apresenta uma leve concentração à esquerda e cauda maior à direita.
Em todas as Grandes Regiões os histogramas também possuem assimetria positiva ( entre 0,57 para o Nordeste e 0,73 para o Centro-Oeste ).
Os Gráficos 3.10, 3.11 e 3.12 apresentam, respectivamente, uma comparação dos resultados em relação à Grande Região do país, à Categoria Administrativa e à Organização Acadêmica, agora levando em conta o desempenho de alunos no componente de Conhecimento Específico da prova.
Pelo Gráfico 3.10, observa -se que existe diferença estatisticamente significativa entre as médias das notas no componente de Conhecimento Específico, das regiões Norte e Centro-Oeste em relação às demais regiões.
A região Nordeste foi a que obteve média mais elevada ( 32,6 ), enquanto a região Centro-Oeste foi a região que obteve média mais baixa ( 29,3 ).
Quanto à Categoria Administrativa ( gráfico 3.11 ), observa -se um comportamento semelhante àquele da parte de Formação Geral e à prova como um todo, ou seja, existe diferença estatisticamente significativa entre as médias das IES Públicas e Privadas, sendo que a maior média foi obtida por alunos de IES Públicas de ensino ( 38,6 ).
Quanto ao Gráfico 3.12, observa -se, mais uma vez, que existe diferença estatisticamente significativa ao nível de 95 \% entre as notas no componente de Conhecimento Específico de Universidades e dos demais tipos de Organização Acadêmica.
Sendo que a média dos Concluintes das Universidades ( 32,9 ) foi maior do que de Centros Universitários ( 31,2 ) e de Faculdades ( 31,1 ). 3.2 ANÁLISE DAS QUESTÕES OBJETIVAS 3.2.1 Componente de Formação Geral A Tabela 3.4 apresenta as estatísticas básicas relativas às oito questões objetivas do componente da prova que abrange a Formação Geral dos estudantes.
A média do Brasil foi 46,0.
A menor média foi encontrada na região Norte ( 43,8 ) e a maior na região Nordeste ( 47,0 ).
As demais médias foram 46,2 na região Sul, 45,9 na região Sudeste e 44,9 na região Centro-Oeste.
O desvio padrão do Brasil foi 20,2, sendo o maior desvio padrão encontrado nas regiões Nordeste e Centro-Oeste ( 20,4 ) e o menor na região Norte ( 19,2 ).
Os demais desvios foram : 20,3 na região Sudeste e 19,9 na região Sul.
As mesmas medianas ( 50,0 ), notas máximas ( 100,0 ) e notas mínimas ( 0,0 ), nas questões objetivas do componente de Formação Geral, foram obtidas em todas as regiões, sem exceção.
A Tabela 3.5 apresenta os índices de facilidade e discriminação ( ponto bisserial ) das questões objetivas do componente de Formação Geral.
Quanto ao índice de facilidade, foram usadas as seguintes cores para diferenciar o nível de dificuldade da questão : • Azul para as questões classificadas com índice muito fácil ( índice \&gt ; =0,86 ), verde para as questões classificadas com índice fácil ( 0,61 a 0,85 ), amarelo para as questões classificadas com médio ( 0,41 a 0,60 ), vermelho para as questões classificadas com difícil ( 0,16 a 0,40 ) e roxo para as questões classificadas com muito difícil ( \&lt ; =0,15 ).
Já quanto ao índice de discriminação, foram usadas as seguintes cores para qualificar a questão : • As questões classificadas com índice fraco receberam a cor roxa ( índice \&lt ; =0,19 ), as classificadas com médio receberam a cor vermelho ( 0,20 a 0,29 ), as classificadas com bom receberam a cor verde ( 0,30 a 0,39 ) e as classificadas com muito bom ( \&gt ; =0,40 ) receberam a cor azul.
As questões objetivas do componente de Formação Geral, segundo o índice de facilidade, foram assim avaliadas : das oito questões, nenhuma teve o índice de facilidade classificado como muito fácil ou fácil.
Seis questões foram tidas como médio, por terem índice de acertos situado na faixa entre 0,41 e 0,60 ( de 41,0 \% a 60,0 \% de acertos ).
Duas questões foram consideradas de dificuldade difícil, situando -se no intervalo entre 0,16 e 0,40 do índice de facilidade, ou seja, houve entre 16,0 \% e 40,0 \% de acertos.
Como já comentado, para análise das questões objetivas relativas à Formação Geral segundo o poder de discriminação, utilizou -se, o índice de discriminação ponto bisserial.
Nesta análise, as questões foram assim avaliadas : seis das oito questões apresentaram índices acima de 0,40 e, assim, foram classificadas com índice muito bom para esse grupo A Tabela 3.5 apresenta os índices de facilidade e discriminação ( ponto bisserial ) das questões objetivas do componente de Formação Geral.
Quanto ao índice de facilidade, foram usadas as seguintes cores para diferenciar o nível de dificuldade da questão : • Azul para as questões classificadas com índice muito fácil ( índice \&gt ; =0,86 ), verde para as questões classificadas com índice fácil ( 0,61 a 0,85 ), amarelo para as questões classificadas com médio ( 0,41 a 0,60 ), vermelho para as questões classificadas com difícil ( 0,16 a 0,40 ) e roxo para as questões classificadas com muito difícil ( \&lt ; =0,15 ).
Já quanto ao índice de discriminação, foram usadas as seguintes cores para qualificar a questão : • As questões classificadas com índice fraco receberam a cor roxa ( índice \&lt ; =0,19 ), as classificadas com médio receberam a cor vermelho ( 0,20 a 0,29 ), as classificadas com bom receberam a cor verde ( 0,30 a 0,39 ) e as classificadas com muito bom ( \&gt ; =0,40 ) receberam a cor azul.
As questões objetivas do componente de Formação Geral, segundo o índice de facilidade, foram assim avaliadas : das oito questões, nenhuma teve o índice de facilidade classificado como muito fácil ou fácil.
Seis questões foram tidas como médio, por terem índice de acertos situado na faixa entre 0,41 e 0,60 ( de 41,0 \% a 60,0 \% de acertos ).
Duas questões foram consideradas de dificuldade difícil, situando -se no intervalo entre 0,16 e 0,40 do índice de facilidade, ou seja, houve entre 16,0 \% e 40,0 \% de acertos.
Como já comentado, para análise das questões objetivas relativas à Formação Geral segundo o poder de discriminação, utilizou -se, o índice de discriminação ponto bisserial.
Nesta análise, as questões foram assim avaliadas : seis das oito questões apresentaram índices acima de 0,40 e, assim, foram classificadas com índice muito bom para esse grupo O Gráfico 3.13, para exemplificar, analisa o comportamento da questão de número 7 de Formação Geral.
Trata -se da terceira questão mais fácil e a que obteve o maior índice de discriminação dessa parte da prova.
Neste gráfico, \textbf{cada} uma das cinco curvas representa o percentual de respostas em determinada alternativa da questão, em função da nota dos estudantes nesta parte da prova ( Formação Geral/Múltipla Escolha ), antes de possíveis eliminações pelo critério do ponto bisserial.
A curva em verde corresponde à alternativa B, a correta para esta questão.
Assim, observa -se que entre os estudantes com notas mais baixas, nessa parte do exame, a situação mais frequente foi a escolha de uma das alternativas incorretas : a alternativa C ( em preto ) ou D ( em roxo ).
Na medida em que a nota aumenta, indicando desempenho melhor nesta parte da prova, aumenta concomitantemente a proporção de estudantes que de alunos ; duas questões tiveram índice de discriminação bom, entre 0,30 e 0,39.
Nenhuma questão foi classificada com um índice fraco e, portanto, nenhuma questão foi eliminada pelo critério ponto bisserial.
O índice de facilidade variou de 0,32 a 0,56, e o de discriminação, de 0,30 a 0,46.
As questões com índice de discriminação muito bom, de números 1, 2, 3, 4, 6 e 7, figuraram entre as mais fáceis desse conjunto : todas classificadas na categoria médio do índice de facilidade.
Em particular, a questão 7 apresentou maior poder discriminatório, com índice 0,46, e foi a terceira mais fácil, com uma proporção de 0,49 de acertos.
A questão de número 5 apresentou índice de facilidade 0,32, ou seja, um quantitativo de 32,0 \% de estudantes conseguiu resolvê -la, dentro do universo de participantes.
Além disso, seu índice de discriminação foi bom, com valor igual a 0,30.
O Gráfico 3.13, para exemplificar, analisa o comportamento da questão de número 7 de Formação Geral.
Trata -se da terceira questão mais fácil e a que obteve o maior índice de discriminação dessa parte da prova.
Neste gráfico, \textbf{cada} uma das cinco curvas representa o percentual de respostas em determinada alternativa da questão, em função da nota dos estudantes nesta parte da prova ( Formação Geral/Múltipla Escolha ), antes de possíveis eliminações pelo critério do ponto bisserial.
A curva em verde corresponde à alternativa B, a correta para esta questão.
Assim, observa -se que entre os estudantes com notas mais baixas, nessa parte do exame, a situação mais frequente foi a escolha de uma das alternativas incorretas : a alternativa C ( em preto ) ou D ( em roxo ).
Na medida em que a nota aumenta, indicando desempenho melhor nesta parte da prova, aumenta concomitantemente a proporção de estudantes que selecionaram a alternativa correta B, atingindo 100 \% para os alunos com 8 acertos.
Essa análise permite verificar como a questão discriminou os grupos de desempenho, justificando o alto índice obtido na questão.
Os gráficos relativos às demais questões de Formação Geral constam do Anexo I. 3.2.2 Componente de Conhecimento Específico A Tabela 3.6 apresenta as estatísticas básicas em relação às questões objetivas do componente de Conhecimento Específico da prova, por Grande Região.
A média do Brasil deste componente foi de 34,9.
A menor média foi observada nas regiões Norte e Centro-Oeste ( 32,4 ) e a maior na região Nordeste ( 36,1 ).
O desvio padrão de todo o Brasil foi 16,6, sendo o menor desvio padrão encontrado na região Norte ( 14,5 ) e o maior na região Nordeste ( 16,9 ).
A mediana de todo o Brasil foi 33,3.
Nas regiões Nordeste, Sudeste e Sul, a mediana foi também 33,3.
Na região Norte e na Centro-Oeste, a mediana foi 28,6.
A nota máxima da prova ( 100,0 ) foi obtida, nas questões objetivas do componente de Conhecimento Específico, por pelo menos um aluno em todas as regiões do Brasil, exceto na Norte ( 90,5 ).
Em todas as regiões, sem exceção, a nota mínima foi zero.
A Tabela 3.7 apresenta os índices de facilidade e discriminação ( ponto bisserial ) das questões objetivas do componente de Conhecimento Específico.
Para facilitar a diferenciação das questões usaram -se as mesmas cores da Tabela 3.5 para as diferentes classificações dos índices de facilidade e de discriminação.
Dentre as questões objetivas da parte da prova relativa ao componente de Conhecimento Específico, nenhuma foi anulada pela Comissão.
Desse modo, a classificação quanto ao índice de facilidade foi estabelecida com base nas 27 questões.
A partir dos índices obtidos, pode -se concluir que a maioria das 27 questões objetivas da prova foi considerada como difícil, 21 ao todo.
Todas as seis questões restantes foram classificadas com índice de facilidade médio.
Já quanto aos índices de discriminação das questões objetivas do componente de Conhecimento Específico da prova, tem -se como resultado a seguinte classificação : dez das 27 questões foram consideradas como boas, enquanto quatro delas tiveram índice de discriminação muito bom.
Assim, para 14 das 27 questões, os índices de discriminação foram bons ou muito bons.
Dentre as demais, sete delas foram classificadas como médias e outras seis como fracas, sendo treze, por conseguinte, a quantidade de questões nos dois patamares mais baixos de discriminação.
Constata -se, assim, que a prova – no que se refere ao componente de Conhecimento Específico – possuía capacidade razoável de discriminar entre aqueles que dominam ou não o conteúdo.
Dentre as questões que alcançaram os maiores índices de discriminação, as de números 13, 15, 18 e 24 classificadas com índice muito bom, situando -se no intervalo de 0,40 a 0,49 do índice, uma delas ( questão 15 ) foi classificada na categoria média e outras três como difícil, quanto ao índice de facilidade.
A questão de número 26 foi a mais difícil dentre as 27 questões específicas, com baixo índice de facilidade, apenas 17,0 \% de acertos.
Essa questão apresentou poder discriminatório igualmente baixo, 0,10, o que comprova ter sido esta a mais difícil para os estudantes.
Destaca -se, também, a questão 25, com índice de facilidade 0,20, o que, em termos percentuais, corresponde a 20,0 \% de estudantes que responderam acertadamente, obtendo, ainda, 0,19 de índice de discriminação.
Tais questões foram, portanto, duas das mais difíceis da prova.
As questões 25 e 26 foram eliminadas do cômputo da nota final pelo critério do ponto bisserial.
Além destas duas, as demais questões com índice fraco de discriminação, questões 9, 10, 22 e 29 também não foram computadas.
A título de exemplo das análises do comportamento das questões objetivas, o Gráfico 3.14 analisa a questão 18 do componente de Conhecimento Específico.
Esta foi uma das questões classificadas como difícil nesta prova, apresentando índice de facilidade 0,34, ou seja, 34,0 \% dos estudantes assinalaram acertadamente a opção E, correspondente ao gabarito.
Seu índice de discriminação foi igual a 0,49, classificado como muito bom.
Neste gráfico, \textbf{cada} uma das cinco curvas representa o percentual de respostas em determinada alternativa da questão 18, em função da nota dos estudantes nesta parte da prova, antes de possíveis eliminações de questões pelo critério do ponto bisserial.
A alternativa correta E, representada no gráfico pela curva em vermelho, foi escolhida em maiores proporções pelos alunos com desempenho melhor nesta parte da prova.
Já as alternativas incorretas, também denominadas distratores, foram selecionadas principalmente por aqueles com notas mais baixas.
A proporção de alunos que selecionou a resposta correta E aumenta gradativamente, chegando a atingir 100 \% para aqueles que acertaram pelo menos 23 questões, enquanto a proporção dos que escolheram alternativas incorretas decai, a partir da primeira nota não nula, como função da nota nesta parte da prova.
Os gráficos relativos às demais questões do componente de Conhecimento Específico constam do Anexo I. 3.3 ANÁLISE DAS QUESTÕES DISCURSIVAS 3.3.1 Componente de Formação Geral As análises dos resultados de desempenho dos estudantes de Administração nas duas questões discursivas relativas à Formação Geral encontram -se na Tabela 3.8 e no Gráfico 3.15.
Na tabela 3.8 observa -se que as notas médias foram mais baixas nesse conjunto de questões do que no das objetivas.
Os estudantes de todo o Brasil obtiveram, em Formação Geral, média 46,0 nas questões objetivas e 38,2 nas questões discursivas.
A mediana também confirma o pior desempenho dos alunos de todo o Brasil nas questões discursivas do componente de Formação Geral.
Enquanto essa estatística foi de 42,5 para questões discursivas, para as questões objetivas a mediana foi de 50,0.
Pode -se, também, notar um aumento do desvio padrão de 20,2, nas questões objetivas do componente de Formação Geral dos alunos de todo o Brasil, para 21,8 nas questões discursivas do mesmo componente.
Como já comentado, a mediana de todo o Brasil, neste componente, foi 42,5, sendo esta encontrada em todas as regiões, exceto na região Centro-Oeste ( 40,0 ).
A nota máxima ( 95,0 ) foi obtida, nas questões discursivas do componente de Formação Geral, por alunos das regiões Nordeste, Sudeste e Sul do Brasil, enquanto na região Norte a nota máxima foi 87,5 e na região Centro-Oeste foi 90,0.
Em todas as regiões a nota mínima foi zero, sem exceção.
O Gráfico 3.15 representa a distribuição das notas nas questões discursivas no componente de Formação Geral.
A moda desta distribuição ocorre no intervalo ( 40 ; 50 ].
Devido à grande quantidade de notas zero e à alta frequência de alunos que deixaram este tipo de questão em branco o intervalo [ 0 ; 10 ] se apresenta como uma segunda moda e com um valor próximo ao da moda principal.
A distribuição possui assimetria à esquerda, coeficiente de assimetria -0,41.
Em todas as regiões o coeficiente de assimetria também é negativo, o maior em módulo, -0,47, na região Norte, e o menor na região Centro-Oeste ( -0,33 ).
Na sequência, os resultados verificados para \textbf{cada} uma das questões discursivas de Formação Geral serão apresentados, estabelecendo -se relações com os conteúdos abordados em \textbf{cada} uma delas.
Os comentários da Banca de docentes corretores a respeito do observado na correção das respostas dos estudantes, suas impressões e conclusões serão apresentados junto à análise de \textbf{cada} questão.
Cumpre esclarecer que, tendo em vista que as questões discursivas de Formação Geral são padronizadas, ou seja, constam de todas as provas, os comentários da Banca são os mesmos para todas as carreiras acadêmicas, sendo direcionados a todos os estudantes que participaram do ENADE/2012.
A seguir, serão analisados os desempenhos da Área de Administração nas duas questões discursivas de Formação Geral do ENADE/2012, comparando os resultados obtidos com comentários para \textbf{cada} questão. 3.3.1.1 Análise da Questão Discursiva 1 do Componente de Formação Geral Os dados de Administração, obtidos a partir das respostas à questão 1, encontram -se na Tabela 3.9 e no Gráfico 3.16.
Nessa questão – de pior desempenho dentre as duas de Formação Geral – os alunos de todo Brasil tiveram média, 34,0.
A maior média para a questão 1 foi obtida na região Nordeste ( 34,5 ), e a menor, na região Centro-Oeste ( 31,7 ).
Quanto à variabilidade das notas, o desvio padrão de todo o Brasil foi 25,2.
O menor desvio padrão foi obtido na região Sul ( 25,0 ) e o maior desvio padrão foi obtido na região Nordeste ( 26,2 ).
A mediana da questão discursiva 1 foi 35,0 para as regiões Norte, Sudeste e Sul, mais alta para a região Nordeste ( 40,0 ) e mais baixa para a Centro-Oeste ( 30,0 ).
Para o Brasil como um todo foi também 35,0.
As notas máximas foram as mesmas para todas as regiões do Brasil ( 95,0 ), sem exceção.
Além disso, a nota mínima foi 0,0 em todas as regiões do país, sem exceção.
O Gráfico 3.16 mostra a distribuição das notas na questão discursiva 1 do componente de Formação Geral.
Observa -se que a maior frequência corresponde aos alunos que deixaram a questão em branco.
Considerando apenas os estudantes que tentaram resolver a questão, a moda ocorre no intervalo ( 30 ; 40 ].
A distribuição das notas possui assimetria positiva ( 0,12 ), assim como em todas as regiões do Brasil ( entre 0,09 para o Nordeste e 0,19 para o Centro-Oeste ). 3.3.1.2 Comentários sobre a correção das respostas à Questão Discursiva 1 A questão aborda um tema bastante polêmico o qual vem sendo amplamente discutido, a conciliação entre o desenvolvimento econômico e a preservação dos recursos naturais.
Busca -se observar a relação entre o estímulo governamental à compra de veículos novos e os impactos ambientais decorrentes desta política pública.
Sua compreensão implica no conhecimento de uma temática atual e na reflexão critica a respeito das decisões governamentais em sua correlação com o cenário econômico internacional e os problemas e soluções de \textbf{ordem} ambiental.
Embora o comando da questão solicitasse explicitamente que a resposta fosse fornecida em um texto único ( “ redija um texto dissertativo sobre... ” ), foi significativamente alta a proporção de respondentes que atendeu a \textbf{cada} item ( a, b e c ) do enunciado de maneira isolada, como partes estanques da mesma resposta.
No que se segue encontra -se uma avaliação da correção de \textbf{cada} um dos três itens. a - “ conceito de desenvolvimento sustentável ” A variedade de respostas a esse item confirmou a importância do tema para avaliar a formação geral dos estudantes e cursos, já que ( i ) é atual ; ( ii ) foi abordado nos meios de comunicação ( especialmente no ano de 2012, em que o Brasil sediou a “ Rio+20 ” ) ; ( iii ) é relevante não apenas no Brasil como em todos os principais fóruns internacionais ; e ( iv ) é comum a diversas áreas de conhecimento.
Outros dois aspectos positivos merecem destaque : ( i ) a grande incidência de respostas que apontaram a relação entre gerações como aspecto essencial da ideia de sustentabilidade socioambiental – no mais, as respostas apresentaram pelo menos a tentativa de indicar uma definição de desenvolvimento sustentável ou a abordagem positiva dessa mesma noção, evidenciado compreensão do enunciado ; e ( ii ) o exercício crítico que a questão propunha, no sentido de exigir dos respondentes a compreensão das contradições de políticas públicas e a apresentação de propostas de soluções para conciliar o incentivo do crescimento econômico e a promoção da sustentabilidade socioambiental, a partir de uma perspectiva transformadora da realidade.
Apesar dos aspectos positivos apontados no parágrafo anterior, a maioria dos estudantes não alcançou o máximo da pontuação.
De forma geral, trataram o desenvolvimento sustentável como sinônimo de não agressão ambiental.
Alguns confundiram desenvolvimento sustentável com desenvolvimento econômico, e outros faziam referência apenas à preocupação com a conservação dos recursos mencionando a preocupação com gerações futuras.
Deve -se, ainda, destacar o significativo número de alunos que apresentou uma definição aparentemente copiada do enunciado de outra questão da prova, como por exemplo, a questão objetiva 31 da prova de Administração. b - “ conflito entre o estímulo à compra de veículos automotores e a promoção da sustentabilidade ” Esse foi o item identificado como aquele no qual as respostas dos estudantes avaliados mais se aproximaram da pontuação máxima prevista no padrão de resposta.
Ou seja, a ampla maioria dos respondentes foi bem-sucedida em apontar a relação entre o aumento da quantidade de veículos automotores em circulação e os danos socioambientais que isso acarreta no cenário urbano, devido ao aumento da emissão de gases poluentes.
No item b, as respostas corretas se concentraram em críticas ao incentivo oferecido pelo governo, referências ao problema dos congestionamentos e ao aumento da poluição.
Contudo, a maioria das respostas não explicitava a correlação entre o aumento da circulação de veículos automotores e o agravamento dos problemas ambientais.
Um número expressivo de alunos apresentou uma resposta em que era dado destaque à política de redução do IPI como um mecanismo para renovação da frota de automóveis ou para a aquisição do transporte próprio e da consequente redução da dependência ao transporte coletivo.
Também foram notados alguns tipos de padrões recorrentes nas tentativas de resposta a este item.
Dois desses padrões se destacaram : ( i ) as declarações genéricas de condenação do “ consumismo ” de uma “ nova classe média ” ou da busca pelo “ status ” que a posse de um automóvel particular proporcionaria, como se fosse possível atribuir a essa \textbf{ordem} de coisas a causa principal do aumento dos congestionamentos nas grandes cidades ; e ( ii ) a repetição da avaliação, presente no enunciado da questão, que aponta de forma superficial a contradição entre o estímulo à demanda por automóveis particulares e a promoção da sustentabilidade. c - “ ações de fomento ao transporte urbano sustentável no Brasil ” Este foi o item que ofereceu maior facilidade de aplicação do padrão de resposta, mas foram poucos os casos de pontuação integral.
Dentre as soluções mais destacadas pelos alunos apareceram : ( i ) o incentivo ao uso de energias alternativas ; e ( ii ) o estímulo ao transporte coletivo ( metrô, trem e ônibus ), além do uso de bicicletas.
Importante destacar o fato de que muitos apresentavam respostas com referências à necessidade de incentivo ao transporte urbano, nas quais este termo era utilizado como sinônimo de transporte coletivo.
A avaliação da questão 1 trouxe à tona problemas estruturais da educação básica, que se refletem no ensino superior.
Foi muito grande o número de provas em que o aluno demonstrou não compreender o enunciado da questão, apresentando respostas que fugiam ao tema proposto, mesmo quando bem redigidas.
Uma resposta recorrente fazia referência às consequências da ampla oferta de crédito à população e aos problemas relacionados à inadimplência, situação que, embora, verdadeira, não \textbf{estava} relacionada à questão proposta.
Houve, ainda, os casos em que as respostas apresentadas não só se afastaram do tema como não reuniam qualquer tipo de argumentação coerente.
E, por fim, destaca -se a existência de respostas incoerentes e desconexas.
Outro aspecto geral diz respeito ao grande número de provas em branco ou em que o aluno apresentou a resposta a outro item no espaço da questão 1.
Considerando -se o objetivo geral deste exame, avaliar a qualidade dos cursos de graduação oferecidos no país e contribuir para a melhoria da qualidade do ensino superior, destaca -se a imensa responsabilidade do processo que envolve a elaboração deste diagnóstico.
Nesse sentido, é urgente alertar para a ocorrência de um número significativo de respostas que se afastam de um padrão mediano de qualidade esperado de um aluno concluinte em um curso de graduação, especialmente em se considerando o fato de que a questão proposta faz referência a tema intensamente explorado pela mídia. 3.3.1.3 Análise da Questão Discursiva 2 do Componente de Formação Geral A Tabela 3.10 mostra que o desempenho dos estudantes na questão 2 ( média 42,4 ) foi superior ao obtido na questão de número 1 ( média 34,0 ).
A região Norte foi aquela onde a média, nessa questão, foi maior ( 45,1 ), e a de menor média foi a região Centro-Oeste ( 40,7 ).
Quanto à variabilidade das notas, o desvio padrão de todo o Brasil foi 24,8, inferior ao obtido na questão de número 1 ( 25,2 ).
O maior desvio nessa questão foi obtido na região Nordeste ( 26,1 ), enquanto o menor foi obtido na região Norte ( 23,6 ).
A mediana foi 50,0 em todas as regiões, exceto na região Centro-Oeste ( 45,0 ).
As notas máximas ( 95,0 ) foram as mesmas em todas as regiões do Brasil, assim como as notas mínimas ( 0,0 ), sem exceção.
O Gráfico 3.17 mostra a distribuição das notas na questão discursiva 2 do componente de Formação Geral.
Neste gráfico, desconsiderando -se a grande quantidade de alunos que deixaram a questão 2 em branco ( em \textbf{torno} de 18 \% ), observa -se a moda no intervalo ( 50 ; 60 ].
O coeficiente de assimetria das notas nesta questão é negativo ( -0,55 ), indicando que a distribuição é assimétrica à esquerda.
O mesmo ocorre para todas as regiões da Brasil ( entre -0,72 na região Norte e -0,48 na Centro-Oeste ). 3.3.1.4 Comentários sobre a correção das respostas à Questão Discursiva 2 O conteúdo da questão era bastante adequado às Diretrizes do INEP para o ENADE, além de atual e de abrangente.
Porém, exigia -se a articulação de muitas variáveis.
Buscando aferir, a partir da apresentação de uma charge e de um texto introdutório, as competências e habilidades relacionadas à leitura, à interpretação, à análise crítica, à capacidade argumentativa, à elaboração de sínteses e de avaliação prospectiva sobre tema da atualidade, a questão tratava da violência e suas repercussões no contexto atual.
O comando da questão solicitava, a partir da definição de violência da OMS e de duas imagens, que o aluno elaborasse um texto, de caráter dissertativo, sobre a “ violência na atualidade ”.
Embora tal solicitação fosse explícita ( “ redija um texto dissertativo sobre... ” ), como na questão 1, foi alta a proporção de alunos que respondeu a \textbf{cada} item do enunciado isoladamente, como partes independentes.
Verificaram -se poucas respostas em branco, como também foram poucas as respostas contendo protestos contra o ENADE, muito frequentes em anos anteriores, o que significa um avanço para a afirmação do ENADE como um exame aceito socialmente. a -"tecnologia e violência"Tanto as charges quanto o texto apresentado no enunciado induziram os alunos a focarem suas respostas nos meios de comunicação.
Além disso, algumas respostas estabeleciam a tecnologia, em si, como causa da violência tanto nas escolas ( em particular ), quanto na sociedade ( no geral ).
Muitos respondentes tiveram dificuldade para fazer a relação entre tecnologia, violência e a escola, de modo adequado, centrando -se na violência nas escolas.
Além disso, a grande maioria das respostas girava em \textbf{torno} da exposição da violência através da TV, internet e jogos digitais, sendo que predominaram as abordagens com ênfase em programas e filmes veiculados pela televisão.
Poucos foram os estudantes que estabeleceram relação entre o contexto socioeconômico e a violência. b -"causas e consequências da violência na escola"Embora a solicitação desse item fosse objetiva e atendida por grande parte dos alunos, inúmeras vezes, observou -se confusão entre o significado de causa e de consequência.
Em outras inúmeras vezes, as respostas apontavam apenas para uma das classes, causas ou consequências.
Foram frequentes as respostas que atribuíam o aumento da violência à ineficiência do poder público, apontando a falta de investimentos, o baixo salário dos professores, a ausência de profissionais qualificados como causas.
Foram frequentes também os que apontavam a violência doméstica, o bullying, e a transferência do imaginário dos filmes e jogos para a realidade, como causas da violência. c -"proposta de solução para o problema da violência na escola"Sem dúvida, há uma grande diversidade de abordagens possíveis para apresentação de uma proposta de solução para o problema da violência, o que o padrão de resposta buscou contemplar.
No entanto, houve grande número de respostas baseadas em opiniões pessoais dos estudantes e vinculadas ao senso comum.
As respostas giraram em \textbf{torno} da necessidade de mais investimentos na educação por parte do poder público e seu rebatimento no maior aparelhamento das escolas, salário dos professores etc.
Além disso, foram muitas as respostas que trataram da necessidade de maior conscientização dos pais, alunos e professores com a realização de um trabalho conjunto que envolvesse todos os segmentos da sociedade.
De um modo geral, os estudantes demonstraram dificuldades em relação ao nível de abrangência exigido, falta de conteúdo, inabilidade na elaboração de texto escrito.
Da grande quantidade de provas corrigidas, o que ficou evidente, foi a dificuldade de estabelecer as conexões requeridas na questão 2.
A violência nas escolas foi bastante associada ao reflexo do meio onde a criança e o adolescente mantém os seus vínculos afetivos.
Surgiram muitas respostas revelando críticas à intolerância no núcleo familiar, à jornada de trabalho dos pais, ao distanciamento da mãe nas atividades domésticas etc.
Em boa parte das respostas, os estudantes lançaram a causa da violência escolar exclusivamente na ausência de participação familiar na vida dos filhos, às vezes criando uma cisão com a proposta da questão, já que abandonavam totalmente a temática da tecnologia em suas reflexões.
Além da ausência dos pais no cotidiano da escola, divórcios e atos de violência contra crianças foram considerados fatores que delegariam para a escola novas funções.
Também foram mencionadas a má formação de professores e a insatisfação com o ambiente escolar.
Em inúmeros casos, ficou evidente o anseio de se fazer da escola um lugar de resgate tanto da criança e do adolescente quanto da família, compreendendo -se a escola em sentido social e assistencial.
Segundo este amplo grupo de concluintes, a escola deveria oferecer -se como centro de cultura e lazer, como espaço para o acompanhamento psicológico e serviço social às comunidades diretamente envolvidas.
Ressalte -se o constante aparecimento das sugestões de horário integral, melhores salários para professores e iniciativas para melhoria na formação de professores.
As propostas de solução para acabar com a violência na escola se detiveram em aspectos muito gerais e poucos alunos detalharam uma proposta de intervenção, apresentando ações concretas.
Dentre as propostas de solução surgiram : educação familiar, ensino religioso, maior segurança nas escolas, com utilização de aparelhos para detecção de armas, e o retorno do papel do professor como um profissional respeitado pelos poderes municipais e estaduais.
Um percentual significativo de provas apresentou resposta simplista, com ênfase em determinados pontos como : religião ( sem demonstração do conceito de laicidade ), políticas sociais ( sem demonstração do domínio de conceitos que envolvem Estado, Governo e Politicas Públicas ) etc.
Foi grande a incidência de respostas que se referiam às leis, à vigilância e maior controle dos pais e professores ; no entanto, não demonstram entender a noção de legalidade para os diferentes segmentos sociais e como os mesmos se encontram submetidos a diferentes formas de poder.
Alguns alunos, ainda que raros, trataram a questão com a profundidade esperada de concluintes de cursos universitários.
Por outro lado, foram vários os textos que apresentaram uma argumentação simplória e desconexa, com forte visão paternalista sobre o papel do governo e do professor, como se tais instâncias pudessem intervir na educação familiar.
Enquanto alguns conseguiram desenvolver o tema proposto com elegância e desenvoltura, encadeando o tema de forma harmônica, com excelente nível de linguagem escrita e correção gramatical, outros sequer conseguiram conectar as ideias e construir um texto razoável, utilizando, inclusive, um linguajar mais propício à linguagem falada do que à escrita.
Cabe também ressaltar a quantidade de respostas carregadas de conteúdo com juízo de valor discriminatório ou extremamente conservador, do tipo : “ mulher não deve trabalhar fora ” ; “ amar e temer a Deus é a solução ” ; obrigatoriedade do ensino religioso e de uma doutrina religiosa ; imposição de uma educação do tipo militar nas escolas ; censura nos meios de comunicação etc.
Outras soluções apontadas para resolver o problema da violência nas escolas \textbf{estavam} associadas à, por exemplo : redução da maioridade penal, defesa da pena de morte, milícias armadas na porta das escolas, punições como o retorno da palmatória.
O Estatuto da Criança e do Adolescente foi uma peça legal criticada, e recebeu, em diversas respostas, o título de responsável pelo sentimento de impunidade de adolescentes infratores.
Contudo, algumas dessas proposições, desde que bem contextualizadas e formuladas com boa argumentação, foram pontuadas de forma a respeitar o direito democrático de pluralidade opinativa.
É imperioso ressaltar que a falta de inteligibilidade segue representando um aspecto merecedor de providências na formação universitária, uma vez que exprime uma habilidade importante para qualquer profissional de nível superior.
Além de terem sido poucas respostas boas e claras, foi possível identificar a deficiência de argumentação, além dos inúmeros problemas gramaticais ( ortografia, concordância, pontuação, regência, etc. ).
Salienta -se, nesse sentido, o uso indiscriminado de letras maiúsculas e minúsculas numa mesma frase, sem obedecer ao critério formal da língua portuguesa. 3.3.2 Componente de Conhecimento Específico Na parte da prova relativa às questões discursivas no componente de Conhecimento Específico ( Tabela 3.11 ), observa -se que a média foi mais baixa do que para as questões discursivas do componente de Formação Geral.
Enquanto no componente de Formação Geral a média para estudantes de Administração de todo o Brasil foi 38,2, na parte de Conhecimento Específico a média foi 14,6.
A maior média deste componente foi obtida pelos estudantes da região Sul ( 15,5 ), e a menor, pelos da região Centro-Oeste ( 11,9 ).
Quanto à variabilidade das notas, o desvio padrão de todo o Brasil foi 17,1.
O maior desvio padrão foi encontrado na região Sudeste ( 17,5 ), e o menor, na região Nordeste ( 15,2 ).
A maior nota máxima, nas questões discursivas do componente de Conhecimento Específico, foi obtida nas regiões Sudeste e Centro-Oeste ( 95,0 ), enquanto a menor nota máxima foi encontrada na região Norte ( 80,0 ).
Além disso, a nota mínima ( 0,0 ) foi obtida por alunos de todas as regiões do Brasil sem exceção.
A mediana do Brasil como um todo foi 8,3, também encontrada na região Norte.
As demais medianas foram : 10,0 nas regiões Sudeste e Sul e 3,3 nas regiões Nordeste e Centro-Oeste.
O Gráfico 3.18 representa a distribuição das notas nas questões discursivas no componente de Conhecimento Específico.
Essa distribuição tem moda no intervalo de [ 0 ; 10 ].
Nota -se que há uma tendência decrescente a partir deste intervalo até o último intervalo.
O coeficiente de assimetria é positivo para todo o Brasil ( 1,08 ) e para todas as regiões ( entre 0,97 para a região Sul e 1,30 para a região Centro-Oeste ), indicando uma assimetria à direita. 3.3.2.1 Análise da Questão Discursiva 3 do Componente de Conhecimento Específico Na questão 3, cujos resultados aferidos encontram -se descritos na Tabela 3.12, a média dos estudantes de todo o Brasil foi 11,5.
A menor média nessa questão foi obtida pelos alunos da região Centro-Oeste ( 8,9 ), enquanto a maior média foi obtida na região Sudeste ( 12,4 ).
Quanto à variabilidade das notas, o desvio padrão de todo o Brasil foi 20,3.
O maior desvio padrão foi obtido na região Sudeste ( 21,3 ), enquanto o menor foi obtido na região Centro-Oeste ( 17,7 ).
A nota máxima, 100,0 pontos, foi alcançada por pelo menos um aluno em todas as regiões, sem exceção.
A mediana e a nota mínima do Brasil como um todo foi zero, o mesmo valor para todas as regiões, sem exceção, indicando que pelo menos metade dos estudantes de \textbf{cada} região recebeu nota zero nesta questão, seja por terem deixado a questão em branco, seja pela inadequação da resposta.
O Gráfico 3.19 mostra a distribuição das notas na questão discursiva 3, do componente de Conhecimento Específico.
Essa distribuição é unimodal com moda nas questões em branco ( também com nota zero ).
O coeficiente de assimetria é positivo para todo o Brasil ( 2,05 ) e para todas as regiões ( entre 1,96 na região Sudeste e 2,45 na região Centro-Oeste ), indicando uma assimetria à direita.
O Gráfico 3.19 mostra a distribuição das notas na questão discursiva 3, do componente de Conhecimento Específico.
Essa distribuição é unimodal com moda nas questões em branco ( também com nota zero ).
O coeficiente de assimetria é positivo para todo o Brasil ( 2,05 ) e para todas as regiões ( entre 1,96 na região Sudeste e 2,45 na região Centro-Oeste ), indicando uma assimetria à direita. 3.3.2.2 Comentários sobre as respostas à Questão Discursiva 3 Muitos estudantes entenderam que a questão deveria ser respondida pela ótica de marketing e/ou finanças, sem utilizar fundamentação teórica para embasar a resposta.
Com isso, foi necessário considerar mais de uma possibilidade alternativa de resposta.
O principal ponto positivo da questão, além do mérito de colocar os estudantes dentro de uma situação existente na prática, foi exigir uma tomada de posição que levava em consideração a análise de aspectos técnicos da situação.
O enunciado da questão apresentava o caso a ser analisado pelos respondentes, com informações dos coeficientes de correlação entre as duas variáveis examinadas nas três cidades.
Assim, para resolução da questão não eram necessários conhecimentos da área financeira ou de marketing, mas sim de estatística ( métodos quantitativos ).
Grande parcela dos estudantes, ao buscar responder o questionamento, demonstrou não saber o que são coeficientes de correlação.
Em relação ao enunciado, podem -se tecer os seguintes comentários sobre \textbf{cada} item : a ) a palavra “ pertinência ” não foi compreendida por muitos alunos ; b ) solicitava duas análises e grande parte dos alunos realizou apenas uma ; c ) pedia uma análise e sugestões.
A maioria realizou apenas a análise.
Assim sendo, o enunciado da questão pareceu não ter sido entendido pela maioria dos alunos, já que desviaram o foco para outros assuntos, como o investimento em projetos comunitários e a discordância dos gerentes.
Além disso, o item a era muito abrangente, pois pedia para comentar a pertinência da correlação, e o aluno, muitas vezes, respondeu apenas “ é ” ou “ não é ” pertinente.
Por não realizarem uma análise de correlação, as respostas dadas pela grande maioria dos estudantes abordaram tomadas de decisão pelo gerente de marketing e/ou pelo gerente financeiro, respostas estas sem fundamentação teórica e que não utilizaram os dados informados, nem as avaliações estatísticas necessárias a uma resposta adequada.
Em relação ao padrão de resposta, pode -se dizer que o mesmo procurou abordar os mais diversos tipos de respostas possíveis, visto que o enunciado era bem amplo, como comentado anteriormente.
Apesar disso, muitos alunos não chegaram sequer ao padrão minimamente requerido.
A média da pontuação mostra o quanto os alunos se distanciaram do padrão esperado.
A correção mostrou, claramente, que os respondentes fizeram diversas interpretações sobre o que \textbf{estava} sendo solicitado.
As respostas, em sua maioria, foram redigidas em texto único, dificultando a pontuação independente dos itens ( a, b e c ), já que a resposta a \textbf{cada} item poderia \textbf{estar} implícita, o que gerou a necessidade de observar a argumentação como um todo.
Além disso, como as correlações apresentadas no enunciado eram baixas para as três cidades, alguns alunos usaram este argumento para justificar suas respostas, corretamente.
Outros, apesar de dizerem que a correlação era baixa, ficaram indecisos e se posicionaram de forma a escolher a cidade C para continuar o investimento.
As duas possibilidades foram consideradas como corretas e pontuadas de acordo.
De maneira geral, com relação aos itens da questão cabem os seguintes comentários : Item a : Muitos alunos, ao invés de responderem sobre a ferramenta estatística utilizada, responderam sobre o uso da área de marketing como ferramenta da administração ou mesmo o uso da responsabilidade social como ferramenta para chamar a atenção dos clientes.
Nesse item, consideramos que os principais problemas foram : a dificuldade de entendimento do significado da palavra “ pertinência ” e a falta de domínio da técnica de análise de correlação.
Item b : A grande maioria se limitou a analisar superficialmente os números.
Poucos deixaram claro o entendimento sobre a análise de correlação.
O mais comum foi dizerem que a correlação era baixa na cidade A, era média na cidade B e que na cidade C a correlação era boa.
Alguns consideraram que a análise geral era boa.
De forma geral, os respondentes achavam o grau de correlação alto nas cidades B e C, bem como no resultado geral.
Alguns alunos fizeram inversão de valores e afirmaram que a cidade A era a que apresentava o melhor índice por ser próximo de zero, mas foram poucas respostas desse tipo.
Algumas respostas apresentavam apenas cópia de dados do enunciado sem fazer nenhuma análise.
Neste item consideramos que os concluintes não compreenderam o que foi \textbf{perguntado} e não compararam o resultado geral com o segmentado por cidade.
Item c : Nesse item, os alunos se preocuparam em defender a área de marketing, afirmando como a mesma era importante, que sempre \textbf{estava} \textbf{certa}, já que pensava no que era melhor para a empresa, enquanto a área de finanças só pensava em economizar.
A tomada de decisão, de fato, foi difícil de ser vista.
Era comum recomendação do tipo : deve haver uma reestruturação na cidade A.
Alguns falaram de novas pesquisas em \textbf{cada} cidade para auxiliar na tomada de decisão.
Neste item destaca -se que as análises, quando apresentadas, eram voltadas para o senso comum, sem qualquer tratamento científico.
Foi constatado que grande parcela dos alunos teve muita dificuldade em realizar a leitura e a interpretação da questão e responderam sem compreender o que se pedia.
É 3.3.2.3 Análise da Questão Discursiva 4 do Componente de Conhecimento Específico A Tabela 3.13 contém as informações relativas à questão 4 do conjunto de questões do componente de Conhecimento Específico.
O desempenho dos estudantes de todo o Brasil nesta questão foi inferior ao desempenho na questão de número 3.
A média geral do Brasil foi 7,7, sendo a menor média registrada na região Centro-Oeste ( 6,3 ) e a maior na região Sudeste ( 8,1 ).
Quanto à variabilidade das notas, o desvio padrão de todo o Brasil foi 13,6.
O maior desvio padrão foi obtido na região Sudeste ( 13,9 ), enquanto o menor foi obtido na região Centro-Oeste ( 12,3 ).
A nota máxima ( 100,0 ) foi atingida em todas as regiões, sem exceção.
A mediana e a nota mínima foram zero para todas as regiões do Brasil, sem exceção.
Ou seja, pelo menos metade dos alunos tirou nota zero nesta questão.
O Gráfico 3.20 representa a distribuição de notas na questão discursiva 4, no componente de Conhecimento Específico.
Como nas outras questões discursivas, o número de estudantes que deixou a questão em branco foi muito grande, correspondendo à moda da distribuição.
O coeficiente de assimetria é positivo para todo o Brasil ( 2,08 ) e para todas as regiões ( entre 2,01 para a região Sudeste e 2,37 para a região Centro-Oeste ), indicando uma assimetria acentuada à direita. 3.3.2.4 Comentários sobre as respostas à Questão Discursiva 4 A questão era, de forma geral, adequada no que diz respeito à linguagem utilizada e à clareza do enunciado.
O padrão de resposta contribuiu para manter estável a avaliação das respostas, apesar do conjunto heterogêneo de respostas redigidas pelos concluintes, e permitiu avaliar e valorizar as mais diferentes respostas.
O desempenho, entretanto, apresentado pelos estudantes na questão 4 foi fraco.
Os estudantes não demonstraram condições mínimas de interpretação, análise, argumentação e aplicação de conceitos.
No geral, mesmo nas poucas respostas consideradas medianas e boas, o estudante apenas citou conceitos.
Foram raríssimos os estudantes que explicaram minimamente os conceitos solicitados pelo enunciado.
Se por um lado cabe destacar o alto número de respostas em branco e o alto número de respostas avaliadas com nota zero, por outro, foram raros os casos de protestos ao ENADE.
Também foram poucos os casos de respostas desconsideradas e os casos de troca de questão.
Poucos alunos responderam ao item ( a ) com exatidão, afirmando que o tipo de estrutura que melhor representa a situação descrita é a estrutura organizacional em rede, que busca integrar a empresa aos stakeholders envolvidos nos seus processos.
Foi observada a predominância de respostas referindo -se a conceitos como organização pós-moderna, estrutura aberta, democrática, estrutura de inovação, modelo participativo.
Tal resultado pode indicar que grande parte dos respondentes não tem conhecimento sobre os tipos de estruturas organizacionais, ancorando suas respostas a características mencionadas como : espírito inovador, internet, organização moderna ( em contraposição com a palavra tradicional ), organização pós-moderna, estrutura inovadora, estrutura horizontal, estrutura democrática, capital intelectual aberto, empresas da nova geração, orgânica, entre outras.
As respostas inadequadas, por sua vez, tratavam de estrutura horizontal, estrutura vertical, estrutura rígida/flexível, sem conexões com aspectos integradores da rede.
Em relação ao item ( b ), as respostas ficaram concentradas em \textbf{torno} da tecnologia de informação ou do sistema de informação, ressaltando a maneira como essas ferramentas agilizam os processos internos das empresas, por permitirem a colaboração de agentes externos para elaboração de novos produtos.
A internet foi largamente mencionada como processo de comunicação para os clientes e outros colaboradores internos e externos das empresas.
Entretanto, as respostas não se referiram à situação apresentada, foi baixo o percentual de alunos que identificou no texto os elementos que se relacionavam com o processo de desenvolvimento.
Por exemplo, várias respostas abordaram a temática apenas a partir da perspectiva da relação cliente/fornecedor.
No item ( c ), a inovação foi citada como a maior motivação para o investimento do fundo de capital de risco, e foi marcante citarem que os colaboradores voluntários seriam clientes ou futuros clientes da empresa.
Foram também dadas respostas que citavam o expressivo volume de vendas, os projetos em fase de acabamento e o prazo de acabamento do produto como incentivo para investimento.
O retorno do capital investido compatível com o risco foi bem menos citado do que as motivações acima descritas.
Cabe ressaltar que várias respostas confundiram ‘ incentivos para investir ’ com ‘ riscos de investimentos ’.
No item ( d ), grande número de respostas demonstrava que os alunos possuem dificuldade em dissociar o setor privado do setor público.
Foi comum confundirem serviços públicos com produtos de empresas.
A colaboração entre cidadão e governo aparecia nas respostas de forma pontual, sem menção à noção de governo aberto.
Poucos conseguiram citar exemplos de soluções inovadoras nos serviços públicos.
Cabe ressaltar que grande parte dos respondentes possui uma concepção de cidadão passivo, entendendo haver apenas a possibilidade de apresentar sugestões aos “ governantes/políticos ” ( em alguns casos tal possibilidade é tratada como se fosse algo a ser criado e/ou oportunizado ).
As respostas focaram apenas práticas empresariais de mercado como a apresentada no texto oferecido pelo enunciado e em como as empresas poderiam trazer benefícios para o governo.
Consideramos que a indicação de que as respostas deveriam ser a respeito da situação apresentada pode ter contribuido para a confusão.
Chamou à atenção a dificuldade dos alunos com a comunicação escrita : foram muitos os erros de português, grafia e concordância, e ficou clara a dificuldade dos estudantes em coordenar as ideias.
Isso se traduziu em respostas fora da norma padrão da língua portuguesa e em respostas fora do contexto solicitado no enunciado da questão.
Perceberam -se, também, dificuldades em sustentação de uma argumentação, mesmo que em poucas linhas, como era o caso do campo previsto para as respostas. 3.3.2.5 Análise da Questão Discursiva 5 do Componente de Conhecimento Específico A Tabela 3.14 contém as informações relativas à questão 5 do conjunto do componente de Conhecimento Específico.
O desempenho dos estudantes nessa questão foi superior ao das questões 3 e 4.
A nota média dos estudantes de todo o Brasil foi 24,5.
A maior média foi registrada na região Sul ( 26,6 ), enquanto a menor média foi registrada na região Centro-Oeste ( 20,4 ).
Quanto à variabilidade das notas, o desvio padrão dos alunos do Brasil, como um todo, foi 29,9.
Enquanto o maior desvio foi encontrado na região Sul ( 30,3 ), o menor foi encontrado na região Centro-Oeste ( 28,3 ).
A nota máxima ( 100,0 ) foi alcançada em todas as regiões do Brasil, sem exceção.
A mediana foi zero para quase todas as regiões do Brasil, exceto para as regiões Sul ( 15,0 ) e Sudeste ( 5,0 ).
A nota mínima foi zero para todas as regiões do Brasil, sem exceção.
O Gráfico 3.21 representa a distribuição das notas da questão discursiva 5 do componente de Conhecimento Específico.
Destaca -se, como nas demais questões discursivas, o grande número de estudantes que deixaram esta questão em branco, correspondendo à moda da distribuição.
Entre as três questões discursivas do Componente de Conhecimento Específico, esta foi a que apresentou um maior espalhamento das notas daqueles que resolveram a questão.
Além da moda na nota zero, aparecem ainda dois máximos locais : ( 20 ; 30 ] e ( 50 ; 60 ].
O coeficiente de assimetria para todos os alunos é 0,84 e para todas as regiões se mantém positivo, variando entre 0,71 na região Sul e 1,09 na região Centro-Oeste. 3.3.2.6 Comentários sobre as respostas à Questão Discursiva 5 Com enunciado claro apresentava -se uma situação verossímil, de fácil entendimento, em linguagem adequada e com bom direcionamento para a resposta dos concluintes.
A questão, com baixo grau de dificuldade, tratava das atribuições básicas de duas das principais áreas funcionais da Administração : recursos humanos e produção.
Em termos de consistência, a contextualização fornecia um interessante pano de fundo para que o estudante refletisse sobre os critérios a serem analisados para uma tomada de decisão estratégica e os explicasse, sendo possível construir um texto dissertativo.
Em uma prova onde o número de questões abertas é limitado, essa característica é muito importante para que as funções da Administração sejam contempladas.
Dessa forma, a questão sobre planejamento apresentava um amplo escopo de possibilidades de respostas.
Tal amplitude pode ser interpretada por alguns como uma dificuldade para avaliar o desempenho dos alunos.
No entanto, o comando da questão era específico e bem direcionado para as áreas funcionais, principalmente no que se refere ao planejamento tático, o que possibilitou a confecção de um padrão de resposta que facilitou a correção.
O enunciado se subdividia em quatro partes : a primeira, em que foi apresentada a contextualização ; a segunda, que resumia a solicitação do diretor de planejamento ; a terceira, que solicitava a redação de um texto dissertativo ; e a quarta que especificava e detalhava o comando ( itens a e b ).
Alguns alunos deram mais atenção à solicitação do diretor do que à resposta aos comandos apresentados nos itens a e b, onde, de fato, encontravam -se as perguntas.
Cabe destacar que as respostas para ambos os itens não demandavam conhecimentos especializados.
Apesar da clareza e simplicidade da questão, muitos concluintes parecem não ter compreendido o que deveriam responder ou negligenciaram alguns aspectos da questão em suas respostas.
Exemplos disso residem no fato da maioria não elaborar um texto dissertativo como solicitado, limitando -se apenas a apresentar uma lista dos critérios críticos das áreas.
Os estudantes, em geral, usaram uma linguagem incoerente com o seu nível de formação e revelaram desconhecimento ou pouca familiaridade com o vocabulário próprio do tema a que se referia à questão.
Muitas respostas pareciam basear -se tão somente no senso comum, sem compromisso com o conhecimento teórico e técnico.
Alguns estudantes chegaram a inventar aspectos que não \textbf{estavam} descritos na contextualização.
O critério de correção estabelecido antecipou a diversidade de respostas dos alunos, e contribuiu para o trabalho de correção, com distribuição de pontuação útil para o processo de correção.
Destaca -se que as situações críticas, observadas na correção amostral, foram incorporadas ao padrão de respostas.
Isso foi fundamental, já que a maioria dos problemas relatados se repetiu nas respostas do universo de estudantes.
Além de outros, a saber : Critérios Críticos da área de Recursos Humanos na decisão de expansão : • oferta de mão de obra especializada na região ; • competências e habilidades para \textbf{cada} cargo ; • recrutamento e seleção ; • necessidade de treinamento ; • custo de mão de obra na região ; • estrutura funcional ( áreas e cargos para nova fábrica ) ; • possibilidade de transferência interna de funcionários. • medidas de incentivo e remuneração variável ; • o fator motivacional como mecanismo de satisfação no trabalho e de retenção ; • conceito de segurança e bem \textbf{estar} no trabalho ; • conceito de comportamento organizacional, clima e cultura.
Uma quantidade insignificante de respostas foi divergente do gabarito.
Critérios Críticos da área de Produção na decisão de expansão :  arranjo físico ;  infraestrutura logística ;  cadeia de fornecimento/suprimento ;  custos de transporte ;  capacidade instalada  demanda e impactos ambientais.
No item b foram poucas as respostas que analisaram : estratégia de produção ; sistemas de produção ; fluxos produtivos ; arranjos produtivos ; organização da produção – produção artesanal, em massa, enxuta ; segurança das operações.
E, de forma geral, quando analisadas, obtiveram uma pontuação muito baixa.
Observou -se também que, com frequência, os respondentes apresentaram os aspectos importantes para a Gerência de Recursos Humanos, deixando de mencionar os critérios relacionados à Produção.
Portanto, as respostas do item a alcançaram nota média superior à do item b.
Os estudantes arriscaram -se mais a escrever sobre recursos humanos do que sobre produção, mesmo sem dominar os conteúdos exigidos.
Adicionalmente, verificou -se um conhecimento superior acerca do vocabulário relativo aos recursos humanos do que ao relativo à produção industrial.
Pode -se inferir que o conteúdo teórico dos cursos de Administração, em particular os das cadeiras ligadas às funções da organização ( marketing, recursos humanos, produção, finanças ), parece dar maior ênfase aos recursos humanos em detrimento da produção.
Isto pode ser analisado sob a ótica de que a produção refere -se, de forma geral, à atividade fim da organização.
Na organização, a produção pode ser industrial ( nos seus mais diferentes processos ), de serviços ou, ainda, comercial.
Essa divisão torna este campo mais complexo para ser ensinado em sala de aula, por dois fatores.
O primeiro, de \textbf{ordem} prática, a distância entre sala de aula e o ambiente prático de produção e, o segundo, de tempo, pela necessidade de abordar os conteúdos com maior profundidade e qualidade, no que se refere à operação ( administração ) industrial, de serviços e comercial.

% matched lemmas: cada, certo, estar, ordem, perguntar, torno
\end{textsample}
